\newif\ifglobal

%%%%%%%%%%%%%%%%%%%%%%%
%%% Compile to view %%%
%%%%%%%%%%%%%%%%%%%%%%%

%%%%%%%%%%%%%%%%%%%%%%%%%%%%%%%%%%%%%%%%%%%%%%%%%%%
%%% Readme:                                     %%%
%%% https://www.overleaf.com/read/bwdjvfjksvjy/ %%%
%%%%%%%%%%%%%%%%%%%%%%%%%%%%%%%%%%%%%%%%%%%%%%%%%%%

% >>> M?: marks a module setting number or important notice

% >>> M4: Set {./relative-path} to external files for this module <<<
\def \modulefiles {./27-SHA}
% To add an external file, use:
% (Replace '---' with actual filename)
% '\input{\modulefiles/tables/---.tex}' for tables
% '\includegraphics{\modulefiles/figures/---}' for figures
% '\subfile{\modulefiles/---}' for register files

% >>> M1: This module is in global project? <<<
% 'YES' - uncomment; 'NO' - comment out
\globaltrue

\ifglobal
    \documentclass[main\_\_CN.tex]{subfiles}

\else
    % Font "FandolSong-Regular" used by ctex does not have GSUB table.
    % It causes warnings that do not affect compilation and can be ignored.
    \PassOptionsToPackage{quiet}{fontspec}

    \documentclass[openany, 10pt]{book}

    % >>> M2: In ./00-shared/preamble-trm-module, also find
    %         settings for visibility of todo notes, watermark <<<
    \usepackage{preamble-trm-module}


    % Enable Chinese typesetting
    \usepackage{ctex}

    % >>> M3: Temporary labels in Chinese
    %         you can add your own labels there <<<
    \myexternaldocument{./00-shared/temp-labels-cn}

\fi %\ifglobal


\setcounter{secnumdepth}{4}

% >>> M5: If this module needs any updates in
% './00-shared/preamble-trm-module.sty'
% (include additional package, etc.), contact your module PM
% or create the file './00-shared/changelog.tex' make the
% necessary updates yourself and document them in this file <<<


\begin{document}


% Sets language into which pre-defined bits of text are translated
\selectlanguage{chinese}

% Do not delete this block
\ifglobal
% Add the following front matter if
% this module is outside of global project
\else
    \InputIfFileExists{readme}{}{}
    {\let\clearpage\relax \listoftodos}
    \tableofcontents
    %\listoftables
    %\listoffigures
\fi


%%%%%%%%%%%%%%%%%%%%%%%%%
%%% TRM Module Begins %%%
%%%%%%%%%%%%%%%%%%%%%%%%%

\hypertarget{sha}{}
\chapter{SHA 加速器 (SHA)} \label{mod:sha}

\section{概述}

\chipname{} 内置 SHA（安全哈希算法）硬件加速器可完成 SHA 运算，具有 \hyperref[sec:sha-typical-sha]{Typical SHA} 和 \hyperref[sec:sha-dma-sha]{DMA-SHA} 两种工作模式。整体而言，相比基于纯软件的 SHA 运算，SHA 硬件加速器能够极大地提高运算速度。

\section{主要特性}

\chipname{} 的 SHA 硬件加速器：

\begin{itemize}
    \item 支持 \href{https://doi.org/10.6028/NIST.FIPS.180-4}{FIPS PUB 180-4 规范}的全部运算标准
    \begin{itemize}
        \item SHA-1 运算
        \item SHA-224 运算
        \item SHA-256 运算
        \item SHA-384 运算
        \item SHA-512 运算
        \item SHA-512/224 运算
        \item SHA-512/256 运算
        \item SHA-512/{\it{t}} 运算
    \end{itemize}
    \item 提供两种工作模式
    \begin{itemize}
        \item Typical SHA 工作模式
        \item DMA-SHA 工作模式
    \end{itemize}
    \item 允许插入 (interleaved) 功能（仅限 Typical SHA 工作模式）
    \item 允许中断功能（仅限 DMA-SHA 工作模式）
\end{itemize}

\section{工作模式简介}
\chipname{} 内置的 SHA 加速器支持两种工作模式。

\begin{itemize}
    \item \hyperref[sec:sha-typical-sha]{Typical SHA 工作模式}：所有数据读写统一通过 CPU 访问完成。
    \item \hyperref[sec:sha-dma-sha]{DMA-SHA 工作模式}：所有读数据通过硬件上的 DMA 完成。具体来说，用户可配置 DMA 控制器，由 DMA 控制器提供 SHA 运算过程中所需的数据信息。因此，可以释放 CPU 执行其他任务。
\end{itemize}

用户可通过配置 \hyperref[regdesc:SHASTARTREG]{SHA\_START\_REG} 或 \hyperref[regdesc:SHADMASTARTREG]{SHA\_DMA\_START\_REG} 选择 SHA 加速器的工作模式，先配置的工作模式生效，具体请见表 \ref{tab:sha-working-mode}。

\begin{longtable}[c]{ |l|l| }
\caption{工作模式选择}
\label{tab:sha-working-mode}
\\\hline
\rowcolor{lightgray}
工作模式      & 选择方式\\
    \hline
    \hyperref[sec:sha-typical-sha]{Typical SHA}      &  \hyperref[regdesc:SHASTARTREG]{SHA\_START\_REG} 置 1   \\
    \hline
    \hyperref[sec:sha-dma-sha]{DMA-SHA}           &  \hyperref[regdesc:SHADMASTARTREG]{SHA\_DMA\_START\_REG} 置 1   \\
    \hline
\end{longtable}

用户可通过配置 \hyperref[regdesc:SHAMODEREG]{SHA\_MODE\_REG} 寄存器选择 SHA 加速器的运算标准，具体请见表 \ref{tab:sha-hash-algo-select}。

\begin{longtable}[c]{|l|c|}
\caption{运算标准选择}
\label{tab:sha-hash-algo-select}
\\\hline
\rowcolor{lightgray}
哈希运算标准 & \hyperref[regdesc:SHAMODEREG]{SHA\_MODE\_REG} 的配置\\
    \hline
    SHA-1 & 0\\
    \hline
    SHA-224 & 1\\
    \hline
    SHA-256 & 2\\
    \hline
    SHA-384 & 3\\
    \hline
    SHA-512 & 4\\
    \hline
    SHA-512/224 & 5\\
    \hline
    SHA-512/256 & 6\\
    \hline
    SHA-512/{\it{t}} & 7\\
    \hline
\end{longtable}

\begin{importantlisting}
\vspace{-0.5em}

    \chipname{} 的 \nameref{mod:digsig} 和 HMAC 模块也会调用 SHA 加速器。此时，用户无法正常访问 SHA 加速器。
\end{importantlisting}

\section{功能描述}

SHA 加速器可以提取信息摘要 (message digest)，其主要工作流程分为两步：\hyperref[sec:sha-preprocessing]{信息预处理}和\hyperref[sec:sha-hash-operation]{哈希运算}。

\subsection{信息预处理}\label{sec:sha-preprocessing}

信息预处理分为三个主要步骤：\hyperref[sec:sha-padding-msg]{附加填充比特}、\hyperref[sec:sha-parsing-msg]{信息解析}和\hyperref[sec:sha-set-initial-hash]{设置初始哈希值}。

\subsubsection{附加填充比特}\label{sec:sha-padding-msg}

SHA 加速器仅能处理长度为 512 位及其整倍数或 1024 位及其整倍数的信息。因此，在将信息送至 SHA 加
速器进行运算前，应先通过软件操作将信息填充为符合要求的长度。

假设待处理信息 \begin{math} M \end{math} 的长度为 \begin{math} m \end{math} 位，则针对不同运算标准的填充步骤见下：

\begin{itemize}
    \item \textbf{SHA-1}、\textbf{SHA-224} 和 \textbf{SHA-256}
        \begin{enumerate}
        \item 首先，在待处理信息后填充 1 个 ``1"；
        \item 随后，再填充 {\it{k}} 个 ``0"。其中，{\it{k}} 为满足 \begin{math}m + 1 + k \ {\equiv} \ 448 \ mod \end{math} 512 的最小非负数解；
        \item 最后，在末尾填充一个 64 位的信息块。该信息块的内容为用二进制表示的待处理信息的长度，即 {\it{m}} 的值。
        \end{enumerate}
    \item \textbf{SHA-384}、\textbf{SHA-512}、\textbf{SHA-512/224}、\textbf{SHA-512/256} 和 \textbf{SHA-512/{\it{t}}}
        \begin{enumerate}
        \item 首先，在待处理信息后填充 1 个 ``1"；
        \item 随后，再填充 {\it{k}} 个 ``0"。其中，{\it{k}} 为满足 \begin{math}m + 1 + k \ {\equiv} \ 896 \ mod \end{math} 1024 的最小非负数解；
        \item 最后，在末尾填充一个 128 位的信息块。该信息块的内容为用二进制表示的待处理信息的长度，即 {\it{m}} 的值。
        \end{enumerate}
\end{itemize}

更多详情，请参考 \href{https://doi.org/10.6028/NIST.FIPS.180-4}{FIPS PUB 180-4 规范} 中的 “5.1 Padding the Message” 章节。

\subsubsection{信息解析}\label{sec:sha-parsing-msg}

在完成信息填充后，我们还需将待处理信息（及其填充）解析为 \begin{math}N\end{math} 个 512 位或 1024 位的信息块。

\begin{itemize}
    \item 对于 \textbf{SHA-1}、\textbf{SHA-224} 和 \textbf{SHA-256}：待处理信息（及其填充）应解析为 \begin{math}N\end{math} 个 512 位的信息块，即 \begin{math}M^{(1)}\end{math}、\begin{math}M^{(2)}\end{math}、 \dots、\begin{math}M^{(N)}\end{math}。一个 512 位信息块包括 16 个 32 位的字 (word)，则第 \begin{math}i\end{math} 个信息块的第一个 32 位字表示为 M$^{(i)}_{0}$，第二个 32 位字表示为 M$^{(i)}_{1}$，\dots，第 16 个 32 位字表示为 M$^{(i)}_{15}$。
    \item 对于 \textbf{SHA-384}、\textbf{SHA-512}、\textbf{SHA-512/224}、\textbf{SHA-512/256} 和 \textbf{SHA-512/\it{t}}：待处理信息（及其填充）应解析为 \begin{math}N\end{math} 个 1024 位的信息块，即 \begin{math}M^{(1)}\end{math}、\begin{math}M^{(2)}\end{math}、 \dots、\begin{math}M^{(N)}\end{math}。一个 1024 位信息块包括 16 个 64 位的字 (word)，则第 \begin{math}i\end{math} 个信息块的第一个 64 位字表示为 M$^{(i)}_{0}$，第二个 64 位字表示为 M$^{(i)}_{1}$，\dots，第 16 个 64 位字表示为 M$^{(i)}_{15}$。
\end{itemize}

SHA 加速器在工作时，每次处理的信息块数据均将按照如下规则写入相应的寄存器中：
\begin{itemize}
    \item \textbf{SHA-1、SHA-224、SHA-256} 将 M$^{(i)}_{0}$ 存放在 \hyperref[regdesc:SHAMNREG]{SHA\_M\_0\_REG} 中，M$^{(i)}_1$ 存放在 \hyperref[regdesc:SHAMNREG]{SHA\_M\_1\_REG}，\dots，M$^{(i)}_{15}$ 存放在 \hyperref[regdesc:SHAMNREG]{SHA\_M\_15\_REG} 中。
    \item \textbf{SHA-384、SHA-512、SHA-512/224、SHA-512/256} 将 M$^{(i)}_0 $ 的高 32 位存放在 \hyperref[regdesc:SHAMNREG]{SHA\_M\_0\_REG} 中，低 32 位存放在 \hyperref[regdesc:SHAMNREG]{SHA\_M\_1\_REG}，M$^{(i)}_1 $ 的高 32 位存放在 \hyperref[regdesc:SHAMNREG]{SHA\_M\_2\_REG} 中，低 32 位存放在 \hyperref[regdesc:SHAMNREG]{SHA\_M\_3\_REG}，...，M$^{(i)}_{15}$ 的高 32 位存放在 \hyperref[regdesc:SHAMNREG]{SHA\_M\_30\_REG} 中，低 32 位存放在 \hyperref[regdesc:SHAMNREG]{SHA\_M\_31\_REG}。

\end{itemize}

\begin{tiplisting}
有关“信息块”及相关概念的描述，请参考 \href{https://doi.org/10.6028/NIST.FIPS.180-4}{FIPS PUB 180-4 规范} 中 “2.1 Glossary of Terms and Acronyms” 章节。
\end{tiplisting}

\subsubsection{哈希初始值 (Initial Hash Value)}\label{sec:sha-set-initial-hash}
在进行哈希运算前，首先必须设置哈希初始值 H$^{(0)}$。不同运算标准的哈希初始值设置要求不同，其中 SHA-1、SHA-224、SHA-256、SHA-384、SHA-512、SHA-512/224、SHA-512/256 等运算的哈希初始值为常量 C，且已经固定在硬件中，无需专门计算。

然而，SHA-512/{\it{t}} 对于不同的 \begin{math}t\end{math} 均需要一个不同的哈希初始值。简单来说，SHA-512/{\it{t}} 是一种基于 SHA-512 的 \begin{math}t\end{math} 位运算标准，其运算结果将截断至 \begin{math}t\end{math} 位。其中，运算标准对 \begin{math}t\end{math} 值的要求为“大于 0 小于 512 且不等于 384 的正整数”。对于不同 {\it{t}} 值的 SHA-512/{\it{t}} 运算，其哈希初始值可通过对 ``SHA-512/t” 字符串的十六进制表示进行 SHA-512 运算获得。不难看出，对于 {\it{t}} 取值不同的 SHA-512/{\it{t}} 运算标准，其不同之处仅在于 {\it{t}} 值不同。

因此，为了简化 SHA-512/{\it{t}} 的哈希初始值计算，我们特别提出了以下方法：

\begin{enumerate}\label{others:sha-512/t}
    \item \textbf{计算 t\_string 和 t\_length}：其中，t\_string 为 {\it{t}} 的字符串信息，长度为 32-bit。t\_length 指明字符串长度信息，长度为 7-bit。根据 t 的取值范围不同，t\_string 和 t\_length 的计算过程如下：

    \begin{itemize}
        \item 如果 $1<=t<=9$，则 $t\_length=7'h48$，$t\_string$ 需要按照如下格式封装:\\
        \begin{longtable}[c]{|l|l|l|}
        \hline
        $8'h30+8'ht_0$ & $1'b1$ & $23'b0$\\
        \hline
        \end{longtable}
        其中，$t_0=t$。

        举例，如果 {\it{t}} = 8，则 $t_0$ = 8，$t\_string$ = $32'h38800000$。

        \item 如果 $10<=t<=99$，则 $t\_length=7'h50$，$t\_string$ 需要按照如下格式封装:\\

        \begin{longtable}[c]{|l|l|l|l|}
        \hline
        $8'h30+8'ht_1$ & $8'h30+8'ht_0$ & $1'b1$ & $15'b0$\\
        \hline
        \end{longtable}
        其中，$t_0=t\%10$，$t_1=t/10$。

        举例，如果 {\it{t}} = 56，则 $t_0$ = 6，$t_1$ = 5，$t\_string$ = $32'h35368000$。

        \item 如果 $100<=t<512$，则 $t\_length=7'h58$，$t\_string$ 需要按照如下格式封装:\\

        \begin{longtable}[c]{|l|l|l|l|l|}
        \hline
        $8'h30+8'ht_2$ & $8'h30+8'ht_1$ & $8'h30+8'ht_0$ & $1'b1$ & $7'b0$\\
        \hline
        \end{longtable}
        其中，$t_0=t\%10$，$t_1=(t/10)\%10$，$t_2=t/100$。

        举例，如果 {\it{t}} = 231，则 $t_0$ = 1，$t_1$ = 3，$t_2$ = 2，$t\_string$ = $32'h32333180$。

    \end{itemize}
    \item \textbf{配置计算哈希初始值所需的寄存器}：用 t\_string 和 t\_length 初始化文本寄存器 \hyperref[regdesc:SHATSTRINGREG]{SHA\_T\_STRING\_REG} 和 \hyperref[regdesc:SHATLENGTHREG]{SHA\_T\_LENGTH\_REG}。

    \item \textbf{计算得到哈希初始值}：对 \hyperref[regdesc:SHAMODEREG]{SHA\_MODE\_REG} 寄存器置 7 选择 SHA-512/{\it{t}} 运算，并对 \hyperref[regdesc:SHASTARTREG]{SHA\_START\_REG} 寄存器置 1，启动 SHA 加速器的运算即可。最后，轮询寄存器 \hyperref[regdesc:SHABUSYREG]{SHA\_BUSY\_REG} 结果为 0，则哈希初始值已计算完毕。
\end{enumerate}


此外，您也可以按照 \href{https://doi.org/10.6028/NIST.FIPS.180-4}{FIPS PUB 180-4 Spec} 中 ``5.3.6 SHA-512/{\it{t}}" 章节的描述计算 SHA-512/{\it{t}} 的哈希初始值，也就是对 “SHA-512/{\it{t}}” 字符串的十六进制表示进行一次“特殊”的 SHA-512 运算，其运算得到的信息摘要即为所需的哈希初始值。这里的“特殊”指本次 SHA-512 运算的哈希初始值为 “SHA-512 运算标准的初始值常量 C 与 0x{}a5 每 8 位进行一次异或位运算后得到的结果”。

\subsection{哈希运算流程}\label{sec:sha-hash-operation}

在完成信息预处理后，\chipname{} SHA 加速器将正式开始哈希运算，最终根据不同运算标准得到不同长度的信息摘要。
正如上文所述，\chipname{} SHA 加速器支持 \hyperref[sec:sha-typical-sha]{Typical SHA} 和 \hyperref[sec:sha-dma-sha]{DMA-SHA} 两种工作模式，下面将对这两种工作模式的具体流程进行介绍。

\subsubsection{Typical SHA 模式下的运算流程}\label{sec:sha-typical-sha}

通常情况下，\chipname{} 的 SHA 会处理完当前信息的所有信息块并生成该信息的信息摘要，之后再开始计算新的信息摘要。不过，\chipname{} SHA 加速器在 Typical SHA 工作模式下还支持 “interleaved” 运算，即在每次运算全部完成前，允许插入其他运算任务。具体来说，在计算完每一个信息块后，用户都可以将存储在 \hyperref[regdesc:SHAHNREG]{SHA\_H\_\regindex{n}\_REG} 寄存器中的信息摘要暂存起来， 然后插入优先级更高的运算任务，包括 Typical SHA 运算和 DMA-SHA 运算。当临时任务结束后，再将之前暂存的信息摘要重新写入 \hyperref[regdesc:SHAHNREG]{SHA\_H\_\regindex{n}\_REG} 中，并继续完成之前中断的计算。

\subparagraph{Typical SHA 的具体运算流程（SHA-512/{\it{t}} 除外）}

    \subfile{\modulefiles/27-SHA-typical-process-except-512t__CN}

\subparagraph{Typical SHA 的具体运算流程（SHA-512/{\it{t}}）}

    \subfile{\modulefiles/27-SHA-typical-process-512t__CN}

\begin{tiplisting}\label{others:sha-note-typical-sha}
\vspace{-2.5em}
\begin{enumerate}
\item 这里，在 SHA 加速器进行硬件运算时，如果存在后续待处理信息块，软件还可以同时将后续信息块写入 \\ \hyperref[regdesc:SHAMNREG]{SHA\_M\_\textcolor{red}{\it{n}}\_REG} 寄存器，以节省时间。
\item 比如重新启动 SHA 加速器完成之前暂停任务的情况。
\item 这里，你可以选择是否需要插入其他任务。如需插入，请前往 \hyperref[others:sha-interleaved]{插入任务工作流程} 具体查看。
\end{enumerate}
\end{tiplisting}

\subfile{\modulefiles/27-SHA-typical-process-interleaved__CN}



\subsubsection{DMA-SHA 模式下的运算流程}\label{sec:sha-dma-sha}

\chipname{} SHA 加速器在 DMA-SHA 工作模式下不支持 “interleaved” 运算方法，即在每次运算全部完成前，不允许插入其他运算任务。这种情况下，用户如有插入运算需求，可将较大信息块进行拆分，并进行多次 DMA-SHA 运算。每次 DMA-SHA 运算之间允许插入其他运算标准的计算任务。

与 Typical SHA 不同，SHA 在 DMA-SHA 工作模式下，运算过程中的数据搬运过程均由硬件完成，具体配置可见章节 \ref{mod:dma} \textit{\nameref{mod:dma}}。

\subfile{\modulefiles/27-SHA-dma-process-except-512t__CN.tex}

\subfile{\modulefiles/27-SHA-dma-process-512t__CN.tex}

\subsection{信息摘要存储}

哈希运算完成之后，计算得到的信息摘要被 SHA 加速器更新至对应的 \hyperref[regdesc:SHAHNREG]{SHA\_H\_\regindex{n}\_REG} (\regindex{n}: 0\verb+~+15) 寄存器中。不同运算标准得到的信息摘要长度也不同，详情见表 \ref{tab:sha-message-digest}：


\begin{table}[H]
    \centering
    \caption{不同运算标准信息摘要的寄存器占用情况}
    \label{tab:sha-message-digest}

    \begin{threeparttable}
    \begin{tabular}{| l | c | l | }
    \hline
        \rowcolor{lightgray}
        哈希运算标准 & 信息摘要长度（位） & 寄存器占用情况\tnote{1} \\\hline

        % table contents
        SHA-1 & 160 & SHA\_H\_0\_REG\verb+ ~ +SHA\_H\_4\_REG\\
        \hline
        SHA-224 & 224 & SHA\_H\_0\_REG\verb+ ~ +SHA\_H\_6\_REG\\
        \hline
        SHA-256 & 256 & SHA\_H\_0\_REG\verb+ ~ +SHA\_H\_7\_REG\\
        \hline
        SHA-384 & 384 & SHA\_H\_0\_REG\verb+ ~ +SHA\_H\_11\_REG\\
        \hline
        SHA-512 & 512 & SHA\_H\_0\_REG\verb+ ~ +SHA\_H\_15\_REG\\
        \hline
        SHA-512/224 & 224 & SHA\_H\_0\_REG\verb+ ~ +SHA\_H\_6\_REG\\
        \hline
        SHA-512/256 & 256 & SHA\_H\_0\_REG\verb+ ~ +SHA\_H\_7\_REG\\
        \hline
        SHA-512/\regindex{t}\tnote{2} & \regindex{t} & SHA\_H\_0\_REG\verb+ ~ +SHA\_H\_\regindex{x}\_REG\\
        \hline
    \end{tabular}

        \begin{tablenotes}
            \item[1] 信息摘要从左至右存放，第一个 word 存放在寄存器 \hyperref[regdesc:SHAHNREG]{SHA\_H\_0\_REG} 中，第二个 word 存放在寄存器 \hyperref[regdesc:SHAHNREG]{SHA\_H\_1\_REG} 中，以此类推。
            \item[2] SHA-512/{\it{t}} 运算标准使用的寄存器与 {\it{t}} 的取值有关。{\it{x}}+1 代表用于存储 {\it{t}} 位信息摘要的 32 位寄存器个数，因此 {\it{x}} = \begin{math}roundup(t/32)\end{math}-1。举例：
                \begin{itemize}
                    \item 当 {\it{t}} = 8 时，则 {\it{x}} = 0，代表最终的信息摘要长度为 8 位，存放在寄存器 \hyperref[regdesc:SHAHNREG]{SHA\_H\_0\_REG} 的高 8 位中；
                    \item 当 {\it{t}} = 32 时，则 {\it{x}} = 0，代表最终的信息摘要长度为 32 位，存放在寄存器 \hyperref[regdesc:SHAHNREG]{SHA\_H\_0\_REG} 中；
                    \item 当 {\it{t}} = 132 时，则 {\it{x}} = 4，代表最终的信息摘要长度为 132 位，存放在寄存器 \hyperref[regdesc:SHAHNREG]{SHA\_H\_0\_REG}、\hyperref[regdesc:SHAHNREG]{SHA\_H\_1\_REG}、\hyperref[regdesc:SHAHNREG]{SHA\_H\_2\_REG}、\hyperref[regdesc:SHAHNREG]{SHA\_H\_3\_REG}，及 \hyperref[regdesc:SHAHNREG]{SHA\_H\_4\_REG} 中。
                    \end{itemize}
        \end{tablenotes}
    \end{threeparttable}
\end{table}

\subsection{ 中断 }\label{sec:sha-Interrupt}

SHA 加速器在 DMA-SHA 工作模式下允许中断发生。
 用户可通过将 \hyperref[regdesc:SHAINTENAREG]{SHA\_INT\_ENA\_REG} 寄存器配置为 1 开启中断。如开启中断功能，SHA 加速器在完成运算时，中断发生。注意，该中断必须由软件将 \hyperref[regdesc:SHAINTCLEARREG]{SHA\_INT\_CLEAR\_REG} 寄存器置为 1 进行清除。由于 SHA 加速器在 Typical SHA 工作模式下的时间开销较小，因此不支持中断功能。

\section{寄存器列表} \label{sec:sha-reg-summ}\hypertarget{sha-reg-summ}{}


本小节的所有地址均为相对于 SHA 加速器基地址的地址偏移量（相对地址），具体基地址请见章节 \ref{mod:sysmem} \textit{\nameref{mod:sysmem}} 中的表 \ref{tab:sysmem-base-address}。

请查看章节 \hyperref[glossary-access-types]{\textit{寄存器的访问类型}}，了解“访问”列缩写的含义。

\subfile{\modulefiles/27-SHA-reg-summ__CN}

\section{ 寄存器 }

本小节的所有地址均为相对于 SHA 加速器基地址的地址偏移量（相对地址），具体基地址请见章节 \ref{mod:sysmem} \textit{\nameref{mod:sysmem}} 中的表 \ref{tab:sysmem-base-address}。

\subfile{\modulefiles/27-SHA-reg__CN}


\end{document}
